% !PRINTED 222
\weberpage{222}
%<<CONT
dass $b = \chi(a)$, $b' = \chi(a')$ die beiden Pole der Substitution
$\chi^{-1}\Theta\chi$ sind. % [sic/mathematics: with b=chi(a), theorem 53.2 requires chi Theta chi^{-1}; the source prints the inverse conjugation order and the ensuing calculation follows that printed order.]
Darin k\"onnen wir $\infty$ und 0 f\"ur $a$ und $a'$ nehmen. % [sic: Weber prints a full stop before the lower-case ``und''; a comma is intended.]
und $\omega$ f\"ur $b$. Dann wird $b'$ ein noch n\"aher zu bestimmender
Pol $\xi$ aus der Reihe (7), und f\"ur $\chi$ erhalten wir die Form

\[
\chi(x) = \frac{\alpha x + \beta}{\gamma x + \delta} = \begin{pmatrix} \alpha, & \beta \\ \gamma, & \delta \end{pmatrix},
\]

worin

\weq{8}{\omega = \frac{\alpha}{\gamma}, \quad \xi = \frac{\beta}{\delta}}

zu setzen ist.
Nun k\"onnen wir f\"ur $\Theta$ jede der Substitutionen $\varepsilon^h x$ nehmen,
wenn nur $h$ nicht durch 5 theilbar ist, und erhalten so die Substitution


\[
\begin{aligned}
\chi^{-1}\Theta\chi &= \begin{pmatrix} \delta, & -\beta \\ -\gamma, & \alpha \end{pmatrix}\begin{pmatrix} \varepsilon^h, & 0 \\ 0, & 1 \end{pmatrix}\begin{pmatrix} \alpha, & \beta \\ \gamma, & \delta \end{pmatrix} \\
&= \begin{pmatrix} \varepsilon^h\alpha\delta - \beta\gamma, & (\varepsilon^h-1)\beta\delta \\ -(\varepsilon^h-1)\alpha\gamma, & \alpha\delta - \varepsilon^h\beta\gamma \end{pmatrix},
\end{aligned}
\]

deren Pole $\omega$ und $\xi$ sein m\"ussen. Nach \S.\,51, (6) m\"ussen also
$\omega$ und $\xi$ die Wurzeln der quadratischen Gleichung

\[
\alpha\gamma x^2 + (\alpha\delta + \beta\gamma)x + \beta\delta = 0
\]

sein, woraus

\[
\frac{\beta\delta}{\alpha\gamma} = \omega\xi, \quad \frac{\delta}{\gamma} + \frac{\beta}{\alpha} = -\omega - \xi
\]

folgt. Aus der zweiten dieser Gleichungen ergiebt sich mittelst (8):

\[
(\alpha + \delta)(\alpha\delta + \beta\gamma) = 0,
\]

und da $\alpha\delta + \beta\gamma$ nicht verschwinden kann, weil sonst $\omega + \xi = 0$
sein m\"usste, was nach (7) nicht m\"oglich ist, so ist $\alpha + \delta = 0$.
Danach erhalten wir f\"ur $\chi$, wenn wir der Einfachheit halber
$\gamma = 1$ annehmen:

\weq{9}{\chi = \begin{pmatrix} \omega, & \beta \\ 1, & -\omega \end{pmatrix}.}

Hierin ist $\beta$, was nach (8) den Werth $-\omega\xi$ hat, noch zu
bestimmen.

Wenn wir die Substitution $\chi$ als bekannt annehmen, so
k\"onnen wir die ganze Ikosa\"edergruppe bilden.

Man hat n\"amlich, wenn man $r$ und $s$ die Zahlen 0, 1, 2, 3, 4
durchlaufen l\"asst, in dieser Gruppe die Substitutionen

\weq{10}{\Theta^r, \quad \psi\Theta^r, \quad \Theta^r\chi\Theta^s, \quad \Theta^r\chi\psi\Theta^s,}
%CONT>>
